% Imporditud: tools/import_eksperimendid.py
% Allikas: Eesti füüsikaolümpiaadi eksperimendi ülesannete
%   kogu (Rael Kalda, 2023), ülesanne Ü19, lahendus L19.
\setAuthor{EFO žürii}
\setRound{piirkonnavoor}
\setYear{2011}
\setNumber{P E1}
\setDifficulty{1}
\setTopic{varia}
\setVahendid{eurosendine münt, 2 joonlauda}
\setPunktid{10}
\setAllikas{Eksperimendi kogu Ü19, lahendus lk 40}
\setLahendusseis{toores_ilma_skeemita}

\prob{Eurosent}
Leidke eurosendise mündi ümbermõõt võimalikult suure täpsusega.

\hint

\solu
Asetame sendi ja joonlaua tasasele pinnale joonisel kujutatud viisil. Sendi pöörame asendisse, kus selle mõni servalähedane märge asub kohakuti ühe joonlaua kriipsuga, mille registreerime alguspunktina. Surume sendi kergelt joonlaudade vahele ja liigutame ühte joonlauda noolega näidatud suunas. Selle tulemusel hakkab sent pöörlema. Libisemise vältimiseks tuleks joonlaudu hoida võimalikult paralleelsena. Peame silmas märget eurosendil ja jätkame liikumist, kuni sent on teinud joonlaua pikkuse piires maksimaalse arvu täispöördeid. Paneme kirja märke lõppkoordinaadi ja pöörete arvu. Koordinaatide vahe jagatud pöörete arvuga annabki eurosendi ümbermõõdu. Kordame katset vähemalt kolm korda ja lõpptulemuseks võtame katsete keskmise. Eurosendi ümbermõõt on 51,1 mm.

\textit{Žürii hindamisskeem on kogumikus lk 40; siia ei ole seda ümber kirjutatud.}
\probend
